% ============================================
% MEM-05
% Gerado de molambudos.md — build_miolo.py
% ============================================
\fragtitle{mem}{M}{MEM}{05}{O Olho Amarelo}
\markboth{MEM-05 --- O Olho Amarelo}{}

\noindent \textit{O que nunca mais fechou}

\noindent \textbf{O que nunca mais fechou}

\vspace{0.5em}\noindent\rule{\textwidth}{0.5pt}\vspace{0.5em}

\noindent Na manhã seguinte, minha mãe me olhou e disse:

\noindent \textit{\textquotedblleft{}Joaquim, teu olho.\textquotedblright{}}

\noindent Eu toquei o rosto. O olho direito estava arregalado. Não conseguia piscar. A pálpebra estava dura, como se tivesse sido costurada para cima.

\noindent Corri para o pote d'água e me inclinei. A água era suja, barrenta, mas dava pra ver o reflexo.

\noindent A íris do olho direito estava amarela.

\noindent Não o amarelo da icterícia, que deixa o olho inteiro cor de gema. Era um amarelo intenso, brilhante, com uma pupila preta que parecia mais funda do que deveria. Como se dentro do meu olho houvesse um poço.

\noindent Minha mãe fez o sinal da cruz.

\noindent \textit{\textquotedblleft{}O que foi que tu viu?\textquotedblright{}}

\noindent Eu não respondi.

\noindent \textit{\textquotedblleft{}Joaquim, pelo amor de Deus, o que foi que tu viu naquela vala?\textquotedblright{}}

\noindent \textit{\textquotedblleft{}Nada\textquotedblright{}\textit{, eu menti. }\textquotedblleft{}Não vi nada.\textquotedblright{}}

\noindent Ela sabia que era mentira. Mas não perguntou de novo. Naquele dia, os soldados vieram buscar as mulheres. Disseram que iam \textquotedblleft{}realocar\textquotedblright{} as moças para o litoral. Minha mãe estava tossindo, tossindo, e mesmo assim a levaram.

\noindent Ela olhou para mim quando o caminhão partiu.

\noindent O olho dela dizia: \textit{não morre.}

\noindent Eu não morri.

\noindent Mas eu queria ter morrido.

\vspace{0.5em}\noindent\rule{\textwidth}{0.5pt}\vspace{0.5em}

\noindent Naquela noite, voltei à vala.

\noindent A criatura estava lá. Dessa vez, ela não estava comendo. Estava deitada no meio dos corpos, olhando para o céu. Quando me viu, sentou.

\noindent \textit{\textquotedblleft{}Sabia que você voltava.\textquotedblright{}}

\noindent Eu fiquei em silêncio.

\noindent \textit{\textquotedblleft{}Vem aqui.\textquotedblright{}}

\noindent Eu fui. Dessa vez, não sentei na borda. Desci. Meus pés afundaram na pilha de corpos. Alguns estavam moles. Outros duros. Um tinha o braço esticado, apontando para o nada. Pisei na mão dele sem querer e ela quebrou com um som seco.

\noindent A criatura tinha um cheiro diferente de perto. Não era só carne. Era sangue grosso, terra molhada, e algo adocicado que lembrava o cheiro de um curral.

\noindent \textit{\textquotedblleft{}Qual é o teu nome?\textquotedblright{}}, ela perguntou.

\noindent \textit{\textquotedblleft{}Joaquim.\textquotedblright{}}

\noindent \textit{\textquotedblleft{}Joaquim.\textquotedblright{}\textit{ Ela repetiu como se estivesse provando a palavra. }\textquotedblleft{}Nome de santo. José, Joaquim, Maria. Nome de quem nasce pra sofrer.\textquotedblright{}}

\noindent \textit{\textquotedblleft{}E o teu?\textquotedblright{}}

\noindent Ela riu. A risada era um som molhado, como água passando por canos sujos.

\noindent \textit{\textquotedblleft{}Não tenho nome. Mas tenho fome. E fome precisa de nome pra chamar.\textquotedblright{}}

\noindent \textit{\textquotedblleft{}Chama de Olho.\textquotedblright{}}

\noindent Ela inclinou a cabeça.

\noindent \textit{\textquotedblleft{}Por quê?\textquotedblright{}}

\noindent \textit{\textquotedblleft{}Porque é a primeira parte de ti que eu vi.\textquotedblright{}}

\noindent Ela sorriu. Os dentes brilharam no escuro.

\noindent \textit{\textquotedblleft{}Então serei Olho. Para ti.\textquotedblright{}}

\noindent Ela estendeu a mão. A mão enorme, de dedos demais, e tocou meu rosto.

\noindent \textit{\textquotedblleft{}Teu olho agora é igual ao meu\textquotedblright{}\textit{, ela disse. }\textquotedblleft{}Tu vês o que eu vejo.\textquotedblright{}}

\noindent \textit{\textquotedblleft{}E o que tu vês?\textquotedblright{}}

\noindent \textit{\textquotedblleft{}Eu vejo a fome.\textquotedblright{}}

\noindent \textit{\textquotedblleft{}Dentro de quem?\textquotedblright{}}

\noindent \textit{\textquotedblleft{}Dentro de todos. Mas a maioria não sabe. Tu vais saber. Tu vais ver em cada rosto o que ele esconde. Vais ver o buraco que cada um tem dentro. E tu vais saber o que preenche esse buraco.\textquotedblright{}}

\noindent \textit{\textquotedblleft{}O quê?\textquotedblright{}}

\noindent Ela inclinou a cabeça para o lado, como um pássaro.

\noindent \textit{\textquotedblleft{}Eu.\textquotedblright{}}

\vspace{0.5em}\noindent\rule{\textwidth}{0.5pt}\vspace{0.5em}

\noindent Na terceira noite, a criatura me mostrou o que eu ia ser.

\noindent Ela tocou minha testa com o dedo --- um dedo longo, de unha preta --- e eu \textit{vi}.

\noindent Vi uma sala. Uma sala grande, cheia de camas. Homens deitados. Homens que gritavam. Homens que não se mexiam. Vi um menino --- mais velho do que eu, mas ainda menino --- sentado num canto. O menino mexia a boca mas não saía som. Ele escrevia num caderno. A letra era minha.

\noindent \textit{\textquotedblleft{}Esse é tu\textquotedblright{}\textit{, disse a criatura. }\textquotedblleft{}Daqui a uns anos.\textquotedblright{}}

\noindent \textit{\textquotedblleft{}Onde é?\textquotedblright{}}

\noindent \textit{\textquotedblleft{}Um lugar que guarda gente como eu. Gente que a sociedade enterrou viva. Tu vai ficar lá muito tempo.\textquotedblright{}}

\noindent \textit{\textquotedblleft{}E tu?\textquotedblright{}}

\noindent \textit{\textquotedblleft{}Eu vou contigo.\textquotedblright{}}

\noindent Ela tocou meu peito. O dedo atravessou a pele, os ossos --- não doeu --- e tocou alguma coisa dentro de mim.

\noindent \textit{\textquotedblleft{}Já estou aqui. Não vou sair mais.\textquotedblright{}}

\noindent \textit{\textquotedblleft{}E se eu não quiser?\textquotedblright{}}

\noindent Ela riu de novo.

\noindent \textit{\textquotedblleft{}Tu não tem escolha, Joaquim. Na noite em que tu sentou na beira da minha vala, tu me convidou. Não com palavras. Com fome. A tua fome me chamou. E eu vim.\textquotedblright{}}

\noindent \textit{\textquotedblleft{}Eu posso ir embora.\textquotedblright{}}

\noindent \textit{\textquotedblleft{}Pode. Mas eu vou dentro de ti.\textquotedblright{}}

\noindent \textit{\textquotedblleft{}E se eu morrer?\textquotedblright{}}

\noindent Ela ficou quieta. Pela primeira vez, o silêncio dela foi diferente. Não era ameaçador. Era... triste.

\noindent \textit{\textquotedblleft{}Se tu morrer, eu espero.\textquotedblright{}}

\noindent \textit{\textquotedblleft{}Espera o quê?\textquotedblright{}}

\noindent \textit{\textquotedblleft{}Até alguém lembrar de ti. E abrir o que tu escreveu. E ler.\textquotedblright{}}

\noindent \textit{\textquotedblleft{}E aí?\textquotedblright{}}

\noindent O olho amarelo dela brilhou no escuro.

\noindent \textit{\textquotedblleft{}Aí eu vou de novo.\textquotedblright{}}

\vspace{0.5em}\noindent\rule{\textwidth}{0.5pt}\vspace{0.5em}

\noindent Quando levantei a cabeça, uma coruja pousada numa árvore seca me observava.

\noindent Não tinha feito barulho. Estava ali, imóvel, me olhando com olhos pretos sem pálpebras. A cabeça dela virou lentamente --- cento e oitenta graus --- e ela fitou a criatura.

\noindent A criatura não olhou para cima. Mas sorriu.

\noindent \textit{\textquotedblleft{}Ela veio te conhecer.\textquotedblright{}}

\noindent \textit{\textquotedblleft{}Quem é?\textquotedblright{}}

\noindent \textit{\textquotedblleft{}É a minha língua no mundo de cima. Tudo que tu disseres, ela leva. Tudo que ela vir, eu sei.\textquotedblright{}}

\noindent A coruja abriu as asas e voou. O vou era silencioso, absorvido pela noite como se o escuro fosse um pano molhado.

\noindent \textit{\textquotedblleft{}Ela vai me seguir?\textquotedblright{}}

\noindent \textit{\textquotedblleft{}Não. Ela vai na frente. Preparando o caminho.\textquotedblright{}}

\vspace{0.5em}\noindent\rule{\textwidth}{0.5pt}\vspace{0.5em}

\noindent Eu tenho 72 anos. Estou escrevendo isto num quarto de hospital psiquiátrico. Os médicos dizem que estou louco. Dizem que a criatura é fruto da minha mente. Que o olho amarelo é uma catarata mal curada.

\noindent Eu sei que não é.

\noindent Ela está aqui agora. Enquanto escrevo. Sinto o peso dela na minha nuca. O hálito doce. Ela está lendo o que escrevo. E está rindo.

\noindent \textit{\textquotedblleft{}Diz pra ele\textquotedblright{}\textit{, ela sussurra. }\textquotedblleft{}Diz que o diário não acaba quando tu morrer.\textquotedblright{}}

\noindent Eu vou morrer em breve. Sei disso. Dr. Oliveira disse que meu coração não aguenta mais.

\noindent Mas o diário vai continuar.

\noindent Não sou eu que escrevo.

\noindent É ela.

\noindent \textit{--- Joaquim Antônio Correia, Paciente 1.260. Último registro.}

\vspace{0.5em}
\noindent\textit{\textbf{↪ Rotas:} \rota{MEM-04} \rotasep \rota{MEM-12} \rotasep \rota{DOC-01}}